\def\probtitle{숭실대입구역}
\def\probno{B} % 문제 번호

\begin{problem}{\probtitle}{표준 입력(stdin)}{표준 출력(stdout)}{1 초}{1024 megabytes}

지하철을 타고 등교하는 도현이는 매번 숭실대입구역에서 내린다. 숭실대입구역에서 출구로 나가기 위해서는 3개의 에스컬레이터를 타야 한다.

길이가 $L$인 에스컬레이터는 가만히 서서 올라가면 $L$초가 걸리고, 체력을 소모하지 않는다. 길이가 $L$인 에스컬레이터를 걸어서 올라가면 $\frac{L}{2}$초 만에 올라갈 수 있지만, $\frac{L}{2}$만큼의 체력을 소모한다. 단, 에스컬레이터를 타는 도중에 올라가는 방법을 바꿀 수는 없다.

도현이는 현재 $H$의 체력을 가지고 있다. 체력을 전부 소모하면 학교에 가서 공부를 할 수 없으므로 남은 체력이 $1$ 이상이 되도록 올라가려고 한다.

3개의 에스컬레이터의 길이 $A$, $B$, $C$와 체력 $H$가 주어졌을 때, 도현이가 체력을 $1$ 이상 남기면서 가장 빠르게 출구로 나갔을 때 걸리는 시간을 출력하는 프로그램을 작성하라.

에스컬레이터의 길이는 항상 짝수로 주어지며, 정답은 정수임이 보장된다.

\InputFile

첫째 줄에 에스컬레이터의 길이 $A$, $B$, $C$가 공백으로 구분되어 주어진다.

둘째 줄에 도현이의 체력 $H$가 주어진다.

\OutputFile

도현이가 체력을 $1$ 이상 남기면서 가장 빠르게 출구로 나갔을 때 걸리는 시간을 출력한다.

\Constraints

\begin{itemize}[topsep=0pt,noitemsep]
    \item 주어지는 모든 수는 정수이다.
    \item $2 \le A, B,C \le 200$
    \item $1 \le H \le 200$
    \item $A,B,C$는 짝수이다.
\end{itemize}

\Example

\begin{example}
    \exmpfile{./example/01.in.txt}{./example/01.out.txt}%
    \exmpfile{./example/02.in.txt}{./example/02.out.txt}%
\end{example}

\Notes

첫 번째 예제에서 길이가 $200$인 에스컬레이터를 걸어서 올라가고, 나머지 에스컬레이터는 가만히 서서 올라가면, 도현이의 체력을 2 남긴 채로 280초 만에 출구로 나갈 수 있다.

\end{problem}